\lstdefinestyle{appendixtxt}{
    basicstyle=\footnotesize\ttfamily,
    showstringspaces=false,
    numbers=left,
    numberstyle=\tiny\color{gray},
    stepnumber=1,
    tabsize=4,
    breaklines=true,
    breakatwhitespace=false,
    frame=lines,
    rulecolor=\color{black}
}

\newcommand{\AppendixPyFile}[2]{%
  \subsection{#1}
  \lstinputlisting[style=appendixpy]{#2}
}

\newcommand{\AppendixPyRangeFile}[4]{%
  \subsection{#1}
  \lstinputlisting[style=appendixpy, firstline=#3, lastline=#4]{#2}
}

\newcommand{\AppendixTxtFile}[2]{%
  \subsection{#1}
  \lstinputlisting[style=appendixtxt]{#2}
}

\chapter{AI工具使用详情}

本附录依据《2026年五一数学建模竞赛人工智能工具使用规定》说明本文中AI工具的使用范围。本文核心建模思路、附件数据处理、参数选择、数值计算结果和论文结论均由参赛队基于题目附件与可复现代码确定；AI工具仅用于部分流程示意图的视觉润色和重新绘制，不用于生成或替换题目附件数据，也不用于决定模型参数和题目交付数值。

\section{工具名称、版本与使用环节}

本文使用的AI工具为ChatGPT，版本/型号为chatgpt-image-2，开发机构为OpenAI，使用日期为2026年5月3日，参考文献按竞赛规定列出为参考文献第1条。

使用环节限定为论文中若干建模流程图和机制示意图的视觉表达优化。被AI辅助润色绘制并在正文采纳的图件包括：各问题关系图、总体建模路线图、数据预处理与特征构造流程图、问题四阶段基线与扰动残差模型结构图、问题五三级预警闭环流程图；另生成问题二复合证据融合示意图作为备选图件，正文最终未启用。上述图件的节点含义、模型关系、文字内容和逻辑方向均由本文建模方案事先确定，AI仅参与图形排版、配色、边框、箭头和文字布局的视觉重绘。

\section{使用原因与人工修改}

初稿曾尝试使用TikZ直接绘制中文流程图，但在编译后的PDF中出现过文字压缩、框内文字溢出、箭头与节点贴合过近、局部元素重叠以及页面比例不协调等视觉偏差。由于这些图件承担的是“解释模型结构和问题关系”的作用，若版式不清晰，容易影响读者对建模证据链的理解。

因此，本文采用“人工确定内容--AI辅助润色--人工复核采纳”的流程：先由参赛队手工确定图件草稿、节点层级、中文说明和模型含义，再使用ChatGPT内置的chatgpt-image-2模型对图件进行学术化视觉润色和重新绘制。生成后逐张进行人工视觉复核，重点检查非白内容是否完整保留、框中文字是否出界、箭头是否压字、元素是否明显重叠、四周留白是否合理。未通过复核的图件不进入论文。

\section{关键交互记录}

图\ref{fig:ai_record_request}和图\ref{fig:ai_record_result}给出本次AI辅助绘图的关键对话记录。对话记录显示，AI使用请求限定在图件重绘和像素级复核，未要求AI生成建模数据、替代计算结果或改写模型结论。

\begin{figure}[H]
\centering
\includegraphics[width=0.96\textwidth]{codex_1.png}
\caption{AI辅助图件重绘的关键请求记录}
\label{fig:ai_record_request}
\end{figure}

\begin{figure}[H]
\centering
\includegraphics[width=0.96\textwidth]{codex_2.png}
\caption{AI辅助图件生成与人工复核结果记录}
\label{fig:ai_record_result}
\end{figure}

\section{采纳范围说明}

正文中凡采用AI辅助润色绘制的图件，均已在图题中标注“AI辅助润色绘制”并引用参考文献第1条。数据驱动图、模型结果表、预测值、异常点清单、阶段节点、预警阈值和附录代码均来自附件数据与可复现程序，不属于AI生成内容。

\chapter{代码复现与数据来源说明}

本附录给出论文主要结果的复现方式、附件数据来源、输出文件对应关系以及核心代码清单。附录中的代码均通过\verb|\lstinputlisting|直接从\texttt{code/}目录引用，避免手工复制造成论文附录与实际运行代码不一致。

\section{运行环境与复现命令}

代码目录为\texttt{code/}，Python依赖由\texttt{uv}管理。进入\texttt{code/}目录后，完整复现顺序如下：

\begin{lstlisting}[style=appendixtxt]
uv run python scripts/run_all.py
uv run python scripts/run_award_audit.py
uv run python scripts/validate_outputs.py
uv run python -m compileall src scripts
uv run pytest
\end{lstlisting}

其中，\texttt{scripts/run\_all.py}生成五问基础求解结果，\texttt{scripts/run\_award\_audit.py}生成模型对照、敏感性分析和预测区间等增强证据，\texttt{scripts/validate\_outputs.py}检查关键表格、图件和约束条件是否满足。正文页数以最终编译结果为准，参考文献与附录不计入正文页数。

\section{原始附件数据来源}

所有模型输入均来自题目提供的五个附件，路径在\texttt{src/slope\_warning/config.py}中统一配置。表\ref{tab:appendix_data_sources}列出附件与问题的对应关系。

\begin{table}[H]
\centering
\small
\caption{附件数据与建模任务对应关系}
\label{tab:appendix_data_sources}
\begin{tabularx}{0.98\textwidth}{L{1.5cm}L{3.0cm}X}
\toprule
问题 & 读取对象 & 建模用途 \\
\midrule
问题一 & 附件1：两组位移时序数据 & A、B两传感器同步位移校准，输出题目表1.1的5个校正值。 \\
问题二 & 附件2：位移时序数据 & 识别三阶段转换节点，建立分阶段位移模型并计算平均速度。 \\
问题三 & 附件3：训练集与实验集 & 训练集缺失补齐、异常识别和贡献排序；实验集表面位移估计。 \\
问题四 & 附件4：训练集与实验集 & 训练集阶段归一化残差建模；实验集指定时刻表面位移预测。 \\
问题五 & 附件5：监测数据 & 枚举六选五变量组合，建立速度分级预警机制。 \\
\bottomrule
\end{tabularx}
\end{table}

需要特别说明的是，附件4和附件5中爆破点距离、单段最大药量的空值表示该时刻无爆破事件，代码在构造爆破冲击特征时将其解释为零强度事件；附件4实验集的表面位移空值为题目要求预测的目标值，不作为普通缺失值补齐。全文没有使用外部生成数据替代附件数据。

\section{主要结果索引}

为保证只阅读PDF即可复核主要数值，附录将关键CSV结果整理为表格。正文引用附录表号，不依赖外部文件路径；表\ref{tab:appendix_outputs}给出附录结果与正文论证的对应关系。

\begin{table}[H]
\centering
\small
\caption{正文关键结果与附录表对应关系}
\label{tab:appendix_outputs}
\begin{tabularx}{0.98\textwidth}{L{1.8cm}L{5.8cm}X}
\toprule
问题 & 复核位置 & 对应正文内容 \\
\midrule
问题一 & 正文表\ref{tab:q1_model_compare}、表\ref{tab:q1_delta_sensitivity}、表\ref{tab:q1_result} & 校准模型对照、Huber阈值敏感性、5个待校正值。 \\
问题二 & 表\ref{tab:appendix_q2_decision}、表\ref{tab:appendix_q2_candidates}、表\ref{tab:appendix_q2_bic} & 复合证据主节点、候选节点解释、分阶段模型阶数选择。 \\
问题三 & 表\ref{tab:appendix_q3_exp_pred} & 实验集表面位移估计节选。 \\
问题四 & 表\ref{tab:appendix_q4_intervals}、表\ref{tab:appendix_q4_shrinkage_cv} & 指定时刻预测值、经验预测区间、残差收缩系数CV对照。 \\
问题五 & 表\ref{tab:appendix_q5_combo_full}、表\ref{tab:appendix_q5_consensus_rank}、表\ref{tab:appendix_q5_warning_summary}、表\ref{tab:appendix_q5_inverse_leadtime} & 六选五组合评价、共识权重敏感性、三级预警事件与逆速度提前量。 \\
\bottomrule
\end{tabularx}
\end{table}

问题二同时保留速度水平跃迁节点作为对照。正文采用复合证据主节点第8108号和第9501号；速度跃迁点第8473号作为第一阶段转换的参考节点，而非正文最终阶段起点。

问题四残差收缩系数的正文表格和附录表\ref{tab:appendix_q4_shrinkage_cv}采用同一组时间块CV结果；$\lambda=0.80$在投影后RMSE上最优，且所有候选值的单调违背数均为0。

\section{公式与代码口径说明}

为避免正文公式与程序实现产生口径差异，表\ref{tab:appendix_formula_code_map}列出关键建模公式及其对应代码位置。正文数值均以这些实现为准。

\begin{table}[H]
\centering
\small
\caption{关键公式与代码实现口径对应表}
\label{tab:appendix_formula_code_map}
\begin{tabularx}{0.98\textwidth}{L{2.4cm}X X}
\toprule
建模环节 & 正文口径 & 对应代码位置 \\
\midrule
爆破冲击 & $I_t=q_t/(d_t+0.5)^2$；非爆破时刻$q_t=0$，并构造6 h、24 h滞后衰减特征 & \path{src/slope_warning/common/features.py}中的\texttt{blast\_impulse}和\texttt{exp\_decay} \\
残差收缩 & 问题四主模型使用$\lambda=0.80$，主验证指标为时间块CV后单调投影RMSE & \path{src/slope_warning/questions/q4_prediction.py}；数值见表\ref{tab:appendix_q4_shrinkage_cv} \\
逆速度窗口 & 最近24 h窗口，即144个10 min采样点；有效样本不少于窗口三分之一，$R^2\ge0.55$且$b<0$ & \path{src/slope_warning/common/diagnostics.py}和\path{src/slope_warning/award/q5.py} \\
三级阈值 & $\theta_1,\theta_2,\theta_3$由6 h平滑速度的阶段分位数、最小等级间隔和持续时间条件共同确定 & \path{src/slope_warning/questions/q5_warning.py}；预警事件见表\ref{tab:appendix_q5_warning_summary} \\
\bottomrule
\end{tabularx}
\end{table}

\chapter{关键结果附表}

本章摘录正文未完全展开但对复核较重要的数值结果。附表保留题目交付、模型选择和预警机制复核所需的关键字段，避免正文依赖外部CSV文件。

\section{问题二复合证据节点输出}

表\ref{tab:appendix_q2_decision}用于明确正文采用的主阶段节点。表\ref{tab:appendix_q2_candidates}用于说明主节点不是单一算法输出，而是由位移趋势、持续加速度和速度跃迁共同给出。

\begin{table}[H]
\centering
\small
\caption{问题二复合证据最终阶段节点}
\label{tab:appendix_q2_decision}
\begin{tabularx}{0.96\textwidth}{L{2.8cm}L{1.5cm}L{3.1cm}L{2.1cm}L{2.1cm}}
\toprule
转换 & 主节点 & 时间 & 位移/mm & 6 h速度/mm/h \\
\midrule
缓慢$\to$加速 & 8108 & 2024-06-29 07:10 & 353.923 & 0.636 \\
加速$\to$快速 & 9501 & 2024-07-08 23:20 & 785.711 & 2.781 \\
\bottomrule
\end{tabularx}
\end{table}

\begin{table}[H]
\centering
\footnotesize
\caption{问题二复合证据候选节点}
\label{tab:appendix_q2_candidates}
\begin{tabularx}{0.98\textwidth}{L{2.1cm}L{2.5cm}L{1.3cm}L{1.2cm}L{2.7cm}X}
\toprule
转换 & 证据类型 & 编号 & 权重 & 时间 & 解释 \\
\midrule
缓慢$\to$加速 & 位移趋势变化 & 5290 & 0.20 & 06-09 17:30 & 分段位移拟合误差最优，反映累计趋势开始偏离。 \\
缓慢$\to$加速 & 持续加速度起点 & 8108 & 0.35 & 06-29 07:10 & 6 h速度连续12 h高于基线阈值，具有提前预警意义。 \\
缓慢$\to$加速 & 速度水平跃迁 & 8473 & 0.45 & 07-01 20:00 & 速度水平分段最优，反映状态显著跃迁。 \\
加速$\to$快速 & 位移趋势变化 & 8871 & 0.25 & 07-04 14:20 & 分段位移拟合误差最优。 \\
加速$\to$快速 & 快速速度持续 & 9567 & 0.25 & 07-09 10:20 & 速度越过加速段与快速段中位水平并持续。 \\
加速$\to$快速 & 速度水平跃迁 & 9501 & 0.50 & 07-08 23:20 & 速度水平分段最优。 \\
\bottomrule
\end{tabularx}
\end{table}

\begin{table}[H]
\centering
\small
\caption{问题二分阶段多项式阶数BIC对照}
\label{tab:appendix_q2_bic}
\begin{tabular}{L{2.4cm}L{2.0cm}L{2.0cm}L{2.0cm}L{1.8cm}}
\toprule
阶段 & 一次BIC & 二次BIC & 三次BIC & 选择阶数 \\
\midrule
缓慢匀速形变 & 52008.8 & -7846.0 & -8847.5 & 三次 \\
加速形变 & 8351.7 & -2064.4 & -2057.9 & 二次 \\
快速形变 & 2394.7 & 1664.1 & 376.8 & 三次 \\
\bottomrule
\end{tabular}
\end{table}

\section{问题三实验集表面位移估计节选}

附件3实验集共有5001行预测结果。表\ref{tab:appendix_q3_exp_pred}摘录前20个编号，用于说明实验集表面位移估计结果已经逐行生成，而非只给出散点图。

\begin{table}[H]
\centering
\footnotesize
\caption{问题三实验集表面位移预测结果节选}
\label{tab:appendix_q3_exp_pred}
\begingroup
\setlength{\tabcolsep}{3pt}
\begin{tabular}{@{}L{0.9cm}L{1.55cm}L{0.9cm}L{1.55cm}L{0.9cm}L{1.55cm}L{0.9cm}L{1.55cm}@{}}
\toprule
编号 & 预测值/mm & 编号 & 预测值/mm & 编号 & 预测值/mm & 编号 & 预测值/mm \\
\midrule
1 & 82.441 & 6 & 79.733 & 11 & 116.909 & 16 & 81.672 \\
2 & 87.621 & 7 & 81.508 & 12 & 88.599 & 17 & 101.727 \\
3 & 84.908 & 8 & 81.508 & 13 & 89.010 & 18 & 103.902 \\
4 & 86.046 & 9 & 81.092 & 14 & 88.158 & 19 & 90.277 \\
5 & 82.302 & 10 & 81.156 & 15 & 89.576 & 20 & 91.829 \\
\bottomrule
\end{tabular}
\endgroup
\end{table}

\section{问题四指定时刻预测区间}

表\ref{tab:appendix_q4_intervals}给出问题四指定时刻预测值、区间和实验集行号。正文表4.1给出同一组结果，本附表保留实验集行号和区间半宽，便于复核指定时刻是否准确命中附件4实验集。

\begin{table}[H]
\centering
\footnotesize
\caption{问题四指定时刻预测值、区间与实验集行号}
\label{tab:appendix_q4_intervals}
\begingroup
\setlength{\tabcolsep}{3pt}
\begin{tabularx}{0.98\textwidth}{@{}L{2.8cm}L{1.2cm}L{0.9cm}L{1.7cm}L{2.8cm}L{1.4cm}@{}}
\toprule
时间点 & 行号 & 阶段 & 预测值/mm & 95\%预测区间/mm & 半宽/mm \\
\midrule
2025-05-09 12:00 & 1125 & 1 & 5.163 & [4.531,\;5.795] & 0.632 \\
2025-05-27 08:00 & 3693 & 2 & 33.594 & [24.085,\;43.102] & 9.509 \\
2025-06-01 12:00 & 4437 & 2 & 61.934 & [52.425,\;71.443] & 9.509 \\
2025-06-03 22:00 & 4785 & 3 & 384.971 & [368.908,\;401.034] & 16.063 \\
2025-06-04 01:40 & 4807 & 3 & 385.482 & [369.419,\;401.546] & 16.063 \\
\bottomrule
\end{tabularx}
\endgroup
\end{table}

\begin{table}[H]
\centering
\footnotesize
\caption{问题四残差收缩系数CV复核表}
\label{tab:appendix_q4_shrinkage_cv}
\begin{tabular}{L{1.5cm}L{2.0cm}L{2.0cm}L{2.1cm}L{2.2cm}}
\toprule
$\lambda$ & 原始RMSE/mm & 原始MAE/mm & 投影后RMSE/mm & 投影后MAE/mm \\
\midrule
0.80 & 2.125 & 1.054 & 2.234 & 1.451 \\
1.00 & 2.142 & 0.860 & 2.470 & 1.548 \\
0.65 & 2.495 & 1.374 & 2.506 & 1.624 \\
0.50 & 3.074 & 1.785 & 3.058 & 1.912 \\
0.35 & 3.766 & 2.233 & 3.754 & 2.269 \\
0.20 & 4.520 & 2.713 & 4.513 & 2.691 \\
0.00 & 5.578 & 3.389 & 5.578 & 3.389 \\
\bottomrule
\end{tabular}
\end{table}

\section{问题五五变量组合完整排序}

正文表7.1展示了主要组合。为满足“枚举全部六选五变量组合”的复核需要，表\ref{tab:appendix_q5_combo_full}列出完整六组结果。保留变量列使用题中变量名称，综合得分越小表示排序越优。

\begin{table}[H]
\centering
\scriptsize
\caption{问题五六种五变量组合完整排序}
\label{tab:appendix_q5_combo_full}
\begingroup
\setlength{\tabcolsep}{3pt}
\begin{tabularx}{0.99\textwidth}{@{}L{5.0cm}L{1.6cm}L{1.8cm}L{2.0cm}L{1.3cm}@{}}
\toprule
保留变量 & 剔除变量 & GBDT位移RMSE/mm & ElasticNet位移RMSE/mm & 综合得分 \\
\midrule
降雨量、孔隙水压力、微震事件数、爆破点距离、单段最大药量 & 干湿入渗系数 & 21.899 & 180.655 & 4.45 \\
降雨量、孔隙水压力、微震事件数、干湿入渗系数、爆破点距离 & 单段最大药量 & 27.399 & 155.954 & 5.35 \\
降雨量、孔隙水压力、微震事件数、干湿入渗系数、单段最大药量 & 爆破点距离 & 29.951 & 155.952 & 5.95 \\
孔隙水压力、微震事件数、干湿入渗系数、爆破点距离、单段最大药量 & 降雨量 & 25.714 & 181.388 & 7.05 \\
降雨量、微震事件数、干湿入渗系数、爆破点距离、单段最大药量 & 孔隙水压力 & 94.221 & 101.589 & 7.20 \\
降雨量、孔隙水压力、干湿入渗系数、爆破点距离、单段最大药量 & 微震事件数 & 215.137 & 247.018 & 12.00 \\
\bottomrule
\end{tabularx}
\endgroup
\end{table}

由表\ref{tab:appendix_q5_combo_full}可见，剔除干湿入渗系数的组合在GBDT位移积分误差上优势明显，虽然ElasticNet对孔隙水压力组合更敏感，但综合非线性拟合和工程解释后，最终选择仍与正文一致。

正文中的共识权重敏感性不重新训练模型，而是直接使用表\ref{tab:appendix_q5_consensus_rank}中的GBDT稳定性名次和ElasticNet稳定性名次，按$J_w(S)=wR_G(S)+(1-w)R_E(S)$代数派生。

\begin{table}[H]
\centering
\footnotesize
\caption{问题五共识排序名次复核表}
\label{tab:appendix_q5_consensus_rank}
\begin{tabularx}{0.98\textwidth}{L{2.0cm}L{2.0cm}L{2.0cm}L{2.0cm}X}
\toprule
剔除变量 & GBDT稳定性名次 & ElasticNet稳定性名次 & 名次差 & 解释 \\
\midrule
干湿入渗系数 & 2.0 & 9.0 & 7.0 & 非线性模型优势明显，需结合阶段非线性拟合解释。 \\
单段最大药量 & 5.0 & 6.0 & 1.0 & 两类模型排序接近，组合结论较稳定。 \\
爆破点距离 & 7.0 & 4.0 & 3.0 & 存在一定排序差异，需结合爆破冲击机理解释。 \\
降雨量 & 6.0 & 9.0 & 3.0 & 存在一定排序差异，降雨效应部分由孔压承接。 \\
孔隙水压力 & 10.0 & 2.0 & 8.0 & 线性模型更敏感，但非线性位移误差明显增大。 \\
微震事件数 & 12.0 & 12.0 & 0.0 & 两类模型均显示剔除后误差较大，应保留。 \\
\bottomrule
\end{tabularx}
\end{table}

\section{问题五预警事件汇总}

表\ref{tab:appendix_q5_warning_summary}给出三级预警事件汇总，正文据此计算首次触发提前量。

\begin{table}[H]
\centering
\small
\caption{问题五三级预警事件汇总}
\label{tab:appendix_q5_warning_summary}
\begingroup
\setlength{\tabcolsep}{3pt}
\begin{tabularx}{0.96\textwidth}{@{}L{2.0cm}L{1.3cm}L{2.1cm}L{2.1cm}L{3.0cm}@{}}
\toprule
等级 & 事件数 & 总持续时间/h & 最大速度/mm/h & 首次触发时间 \\
\midrule
关注 & 7 & 322.167 & 10.115 & 2024-01-17 04:00 \\
预警 & 2 & 184.167 & 10.115 & 2024-02-01 08:10 \\
严重预警 & 9 & 110.500 & 10.115 & 2024-02-01 12:20 \\
\bottomrule
\end{tabularx}
\endgroup
\end{table}

正文中的预警提前量表由表\ref{tab:appendix_q5_warning_summary}的首次触发时间与第一阶段转换时刻相减得到；逆速度提前量统计见表\ref{tab:appendix_q5_inverse_leadtime}，未重新构造预警事件。

\begin{table}[H]
\centering
\small
\caption{问题五逆速度提前量统计}
\label{tab:appendix_q5_inverse_leadtime}
\begin{tabular}{L{1.4cm}L{1.4cm}L{2.4cm}L{2.4cm}L{2.4cm}}
\toprule
阶段 & 记录数 & 提前量中位数/h & 最小提前量/h & 最大提前量/h \\
\midrule
1 & 6 & 12.96 & 11.42 & 14.98 \\
2 & 25 & 5.40 & 0.21 & 21.17 \\
3 & 3 & 2.83 & 1.12 & 4.56 \\
\bottomrule
\end{tabular}
\end{table}

\chapter{核心代码摘录}

本章先列出与模型复现和关键算法直接相关的核心代码片段，便于读者快速定位五个问题的主要求解逻辑。后续附录章节进一步给出完整源程序文件，包括主求解脚本、公共模块、增强审查模块、图件生成脚本、配置文件和测试入口，以满足PDF内复核代码完整性的要求。

\section{复现入口与输出校验}

五个问题由统一入口顺序运行，主结果包括表格、模型摘要和图件。输出校验脚本检查关键表图、无缺失、非负约束、问题四单调性以及问题五阈值递增性。

\AppendixPyRangeFile{五问统一求解入口}{code/scripts/run_all.py}{18}{32}
\AppendixPyRangeFile{关键输出一致性校验}{code/scripts/validate_outputs.py}{72}{116}

\section{问题一：Huber稳健仿射校准}

问题一的核心在于用IRLS求解Huber稳健仿射模型，降低异常同步观测点对传感器校准的影响。

\AppendixPyRangeFile{Huber-IRLS校准函数}{code/src/slope_warning/questions/q1_calibration.py}{14}{39}

\section{问题二：速度主导阶段识别}

问题二以Hampel清洗后的6 h滚动中位速度作为主证据，再用两断点分段和位移分段拟合交叉验证阶段节点。

\AppendixPyRangeFile{速度清洗、两断点识别与阶段拟合}{code/src/slope_warning/questions/q2_segmentation.py}{25}{38}
\AppendixPyRangeFile{阶段转换判别准则输出}{code/src/slope_warning/questions/q2_segmentation.py}{71}{89}

\section{问题三：连续状态量与稀疏驱动量分治}

问题三将孔压、深部位移、表面位移作为连续状态量补齐和诊断，将降雨量、微震事件数作为稀疏驱动量保留真实峰值，再用分组置换重要性解释变量贡献。

\AppendixPyRangeFile{多源变量预处理与异常标记}{code/src/slope_warning/questions/q3_fusion.py}{55}{108}
\AppendixPyRangeFile{分组置换重要性计算}{code/src/slope_warning/questions/q3_fusion.py}{139}{153}

\section{问题四：阶段归一化基线与残差收缩}

问题四先对每个阶段建立单调PCHIP基线，再用扰动残差模型修正局部偏差，并通过时间块交叉验证选择残差收缩系数，最后进行单调投影。

\AppendixPyRangeFile{阶段归一化时间与PCHIP基线}{code/src/slope_warning/questions/q4_prediction.py}{31}{78}
\AppendixPyRangeFile{单调投影和残差收缩交叉验证}{code/src/slope_warning/questions/q4_prediction.py}{122}{177}

\section{问题五：变量组合筛选与三级预警}

问题五用全部六选五组合进行模型比较，并在速度序列上叠加持续时间条件形成三级预警事件。

\AppendixPyRangeFile{五变量组合时间块验证}{code/src/slope_warning/questions/q5_warning.py}{53}{79}
\AppendixPyRangeFile{三级阈值与持续事件提取}{code/src/slope_warning/questions/q5_warning.py}{81}{173}

\chapter{分问题主求解脚本}

本章列出问题一至问题五的主求解脚本。与前一章的核心片段不同，本章保留完整执行逻辑，便于复核每一问从附件读取、模型求解到结果表和图件输出的全过程。

\section{问题一：稳健传感器校准}

\AppendixPyFile{问题一主脚本 \texttt{q1\_calibration.py}}{code/src/slope_warning/questions/q1_calibration.py}

\section{问题二：速度主导阶段识别}

\AppendixPyFile{问题二主脚本 \texttt{q2\_segmentation.py}}{code/src/slope_warning/questions/q2_segmentation.py}

\section{问题三：多源补齐、异常识别与贡献分析}

\AppendixPyFile{问题三主脚本 \texttt{q3\_fusion.py}}{code/src/slope_warning/questions/q3_fusion.py}

\section{问题四：阶段归一化残差预测}

\AppendixPyFile{问题四主脚本 \texttt{q4\_prediction.py}}{code/src/slope_warning/questions/q4_prediction.py}

\section{问题五：变量筛选与速度预警}

\AppendixPyFile{问题五主脚本 \texttt{q5\_warning.py}}{code/src/slope_warning/questions/q5_warning.py}

\chapter{公共预处理、分段与特征模块}

本章列出被多个问题复用的公共模块。保留这些模块可以说明全文并非每问各自临时处理数据，而是在统一时间索引、统一滚动窗口、统一评价指标和统一特征构造框架下完成建模。

\section{预处理模块}

\AppendixPyFile{缺失补齐、Hampel规则与滚动窗口 \texttt{preprocessing.py}}{code/src/slope_warning/common/preprocessing.py}

\section{分段与阶段拟合模块}

\AppendixPyFile{两断点搜索与分阶段拟合 \texttt{segmentation.py}}{code/src/slope_warning/common/segmentation.py}

\section{特征构造模块}

\AppendixPyFile{窗口特征、扰动特征与六选五组合 \texttt{features.py}}{code/src/slope_warning/common/features.py}

\section{指标与诊断模块}

\AppendixPyFile{评价指标与时间块划分 \texttt{metrics.py}}{code/src/slope_warning/common/metrics.py}
\AppendixPyFile{单调性与残差诊断 \texttt{diagnostics.py}}{code/src/slope_warning/common/diagnostics.py}

\section{数据读取与图件公共模块}

\AppendixPyFile{附件读取与输出目录管理 \texttt{io.py}}{code/src/slope_warning/common/io.py}
\AppendixPyFile{论文图件绘制公共函数 \texttt{plotting.py}}{code/src/slope_warning/common/plotting.py}

\chapter{配置、运行入口与测试脚本}

本章列出工程配置、统一运行入口、输出校验、图件生成和测试脚本。至此，论文主结果、增强审查结果和图件均可由PDF中列出的代码复现；依赖环境由\texttt{pyproject.toml}说明。

\section{依赖配置与路径配置}

\AppendixTxtFile{Python依赖配置 \texttt{pyproject.toml}}{code/pyproject.toml}
\AppendixPyFile{附件路径与输出目录配置 \texttt{config.py}}{code/src/slope_warning/config.py}
\AppendixPyFile{包初始化文件 \texttt{\_\_init\_\_.py}}{code/src/slope_warning/__init__.py}
\AppendixPyFile{问题模块初始化文件 \texttt{questions/\_\_init\_\_.py}}{code/src/slope_warning/questions/__init__.py}
\AppendixPyFile{公共模块初始化文件 \texttt{common/\_\_init\_\_.py}}{code/src/slope_warning/common/__init__.py}

\section{运行入口与输出校验}

\AppendixPyFile{五问统一运行脚本 \texttt{run\_all.py}}{code/scripts/run_all.py}
\AppendixPyFile{特等奖增强审查入口 \texttt{run\_award\_audit.py}}{code/scripts/run_award_audit.py}
\AppendixPyFile{稳健性审查入口 \texttt{run\_robustness\_audit.py}}{code/scripts/run_robustness_audit.py}
\AppendixPyFile{输出一致性校验脚本 \texttt{validate\_outputs.py}}{code/scripts/validate_outputs.py}
\AppendixPyFile{论文图件生成脚本 \texttt{make\_paper\_diagrams.py}}{code/scripts/make_paper_diagrams.py}

\section{单元测试}

\AppendixPyFile{公共流程单元测试 \texttt{test\_common\_pipeline.py}}{code/tests/test_common_pipeline.py}

\chapter{增强审查完整代码}

本章列出模型质量审查与稳健性分析所用的完整代码。增强审查不改变题目主结果，而是生成模型对照、bootstrap区间、敏感性分析、预测区间、变量组合稳定性和逆速度提前量等证据，用于支撑正文中的模型评价和附录结果表。

\section{增强审查入口与兼容封装}

\AppendixPyFile{增强审查统一封装 \texttt{award\_audit.py}}{code/src/slope_warning/award_audit.py}
\AppendixPyFile{增强审查包初始化 \texttt{award/\_\_init\_\_.py}}{code/src/slope_warning/award/__init__.py}

\section{增强审查公共函数}

\AppendixPyFile{增强审查公共工具 \texttt{award/common.py}}{code/src/slope_warning/award/common.py}

\section{分问题增强审查}

\AppendixPyFile{问题一模型对照与bootstrap审查 \texttt{award/q1.py}}{code/src/slope_warning/award/q1.py}
\AppendixPyFile{问题二复合证据与节点敏感性审查 \texttt{award/q2.py}}{code/src/slope_warning/award/q2.py}
\AppendixPyFile{问题三异常阈值与变量贡献审查 \texttt{award/q3.py}}{code/src/slope_warning/award/q3.py}
\AppendixPyFile{问题四消融、收缩与预测区间审查 \texttt{award/q4.py}}{code/src/slope_warning/award/q4.py}
\AppendixPyFile{问题五变量筛选与预警提前量审查 \texttt{award/q5.py}}{code/src/slope_warning/award/q5.py}


